% Sample table integration for CAMELS-RU manuscript
% This file shows how to include the tables in your sections

% In section 02_study_area.tex or 03_data_sources_methods.tex:
% Add after describing the data:

The CAMELS-RU dataset comprises \ntotal watersheds with comprehensive hydrological monitoring 
(Table~\ref{tab:data_coverage}). The catchments span a wide range of sizes 
(Table~\ref{tab:watershed_size}), with median area of \medianarea, ensuring representation 
of diverse hydrological processes from headwater streams to major river systems.

\begin{table}[htbp]
\centering
\caption{Summary of hydrological data coverage for CAMELS-RU dataset (2008--2023).}
\label{tab:data_coverage}
\begin{tabular}{lrr}
\toprule
\textbf{Category} & \textbf{Count} & \textbf{\%} \\
\midrule
\multicolumn{3}{l}{\textit{Discharge Data}} \\
\quad Total discharge gauges & 2183 & 82.8 \\
\quad Full records (decent quality) & 1021 & 46.8 \\
\quad Full records (poor quality) & 76 & 3.5 \\
\quad Partial records (decent quality) & 769 & 35.2 \\
\quad Partial records (poor quality) & 317 & 14.5 \\
\midrule
\multicolumn{3}{l}{\textit{Water Level Data}} \\
\quad Total level gauges & 2687 & 101.9 \\
\quad Full records (decent quality) & 1358 & 50.5 \\
\quad Full records (poor quality) & 20 & 0.7 \\
\quad Partial records (decent quality) & 861 & 32.0 \\
\quad Partial records (poor quality) & 82 & 3.1 \\
\quad Freezing-affected & 140 & 5.2 \\
\quad Negative values detected & 226 & 8.4 \\
\midrule
\multicolumn{3}{l}{\textit{Co-located Measurements}} \\
\quad Gauges with both Q and H & 2031 & 77.0 \\
\quad Discharge only & 152 & 5.8 \\
\quad Water level only & 656 & 24.9 \\
\bottomrule
\end{tabular}
\begin{tablenotes}
\small
\item Note: Percentages for discharge and water level data are relative to their respective totals. Co-located percentages are relative to total watersheds (2638).
\item Full records: $>$80\% completeness during study period; Partial: 10--80\% completeness.
\item Decent quality: operational sensors with regular maintenance; Poor quality: gaps, sensor issues, or maintenance concerns.
\end{tablenotes}
\end{table}

\begin{table}[htbp]
\centering
\caption{Watershed size distribution in CAMELS-RU dataset.}
\label{tab:watershed_size}
\begin{tabular}{lrr}
\toprule
\textbf{Size Category} & \textbf{Count} & \textbf{\%} \\
\midrule
$<$ 100 km$^2$ (headwater) & 137 & 5.2 \\
100--2\,000 km$^2$ (small) & 1026 & 38.9 \\
2\,000--10\,000 km$^2$ (medium) & 641 & 24.3 \\
10\,000--50\,000 km$^2$ (large) & 453 & 17.2 \\
50\,000--200\,000 km$^2$ (very large) & 198 & 7.5 \\
$>$ 200\,000 km$^2$ (major rivers) & 183 & 6.9 \\
\midrule
\textbf{Total} & \textbf{2638} & \textbf{100.0} \\
\bottomrule
\end{tabular}
\begin{tablenotes}
\small
\item Median watershed area: 2923 km$^2$. Size categories follow established hydrological basin classification.
\item Detailed analysis focuses on watersheds $<$ 50\,000 km$^2$ (n = 2251, 85.3\%) to minimize routing effects and human regulation.
\end{tablenotes}
\end{table}


% In section 03_data_sources_methods.tex:
% Add after discussing meteorological forcing:

Three precipitation datasets were evaluated for the CAMELS-RU domain (Table~\ref{tab:precipitation_comparison}). 
ERA5-Land shows the highest mean annual precipitation (\erafiveannual) and strongest correlation with 
observed discharge ($r = 0.087$), while MSWEP and GPCP provide lower estimates with similar spatial patterns.

\begin{table}[htbp]
\centering
\caption{Precipitation dataset comparison for CAMELS-RU catchments (2008--2023).}
\label{tab:precipitation_comparison}
\begin{tabular}{lrrrr}
\toprule
\textbf{Dataset} & \textbf{N} & \textbf{Mean P} & \textbf{Std P} & \textbf{P--Q Corr} \\
 & \textbf{Gauges} & \textbf{(mm/yr)} & \textbf{(mm/yr)} & \textbf{(lag 0)} \\
\midrule
ERA5-Land & 1721 & 893 & 246 & 0.087 \\
MSWEP v2.8 & 1721 & 604 & 183 & 0.066 \\
GPCP v3.2 & 1721 & 653 & 167 & 0.081 \\
\bottomrule
\end{tabular}
\begin{tablenotes}
\small
\item P--Q Corr: Pearson correlation between annual precipitation sum and annual discharge at lag 0.
\item N = 1{,}721: catchments with concurrent discharge and complete precipitation coverage from all three products.
\item ERA5-Land shows highest mean precipitation and strongest correlation with discharge.
\item All datasets show similar spatial patterns but differ in absolute magnitudes (see Section~\ref{sec:technical_validation}).
\item GPCP coverage: 2008--2021 only.
\end{tablenotes}
\end{table}


% In section 04_attributes_signatures.tex:
% Add after describing static attributes:

CAMELS-RU includes \nattributes static catchment attributes derived from HydroATLAS v1.0 
(Table~\ref{tab:hydroatlas_attributes}), spanning six thematic categories: topography and climate, 
cryosphere, land cover, soil properties, hydrogeology, and flood regulation. These attributes 
capture the physical and physiographic diversity of Russian catchments.

\begin{table*}[htbp]
\centering
\caption{Static catchment attributes derived from HydroATLAS v1.0 included in CAMELS-RU (\nattributes primary subset). Initial selection retained attributes matching upstream and sub-basin aggregation conventions; 12 inter-correlated features were subsequently removed (see Sect.~\ref{sec:data_sources_methods}). The full set of 288 HydroATLAS attributes is provided in the released CSV.}
\label{tab:hydroatlas_attributes}
\small
\begin{threeparttable}
\begin{tabular}{lllp{7cm}}
\toprule
\textbf{Category} & \textbf{Attribute} & \textbf{Unit} & \textbf{Description} \\
\midrule
\multirow{5}{*}{Climate \& water balance}
& ele\_mt\_uav & m & Mean elevation \\
& pre\_mm\_uyr & mm/yr & Mean annual precipitation \\
& pet\_mm\_uyr & mm/yr & Mean annual potential evapotranspiration \\
& aet\_mm\_uyr & mm/yr & Mean annual actual evapotranspiration \\
& snw\_pc\_uyr & \% & Mean annual snow cover extent \\
\midrule
\multirow{5}{*}{Land cover}
& for\_pc\_use & \% & Forest cover \\
& crp\_pc\_use & \% & Cropland cover \\
& pst\_pc\_use & \% & Pasture/grassland cover \\
& ire\_pc\_use & \% & Irrigated area \\
& pac\_pc\_use & \% & Protected area coverage \\
\midrule
\multirow{2}{*}{Cryosphere}
& gla\_pc\_use & \% & Glacier cover \\
& prm\_pc\_use & \% & Permafrost extent \\
\midrule
\multirow{3}{*}{Soil}
& cly\_pc\_uav & \% & Clay content (mean) \\
& slt\_pc\_uav & \% & Silt content (mean) \\
& snd\_pc\_uav & \% & Sand content (mean) \\
\midrule
\multirow{2}{*}{Hydrogeology}
& gwt\_cm\_sav & cm & Groundwater table depth \\
& kar\_pc\_use & \% & Karst area extent \\
\midrule
\multirow{2}{*}{Water bodies}
& inu\_pc\_ult & \% & Inundation extent \\
& lka\_pc\_use & \% & Lake/reservoir area \\
\midrule
\multirow{3}{*}{Human impact}
& ppd\_pk\_uav & pers/km$^2$ & Population density \\
& urb\_pc\_use & \% & Urban/built-up area \\
& gdp\_ud\_sav & USD/km$^2$ & GDP density \\
\bottomrule
\end{tabular}
\begin{tablenotes}
\small
\item All attributes are watershed-level aggregates computed individually for each manually verified catchment boundary using area-weighted spatial averaging.
\item Attribute codes follow HydroATLAS naming: suffix indicates aggregation (uav = upstream area-weighted average, use = upstream spatial extent, uyr = upstream annual mean, sav = sub-basin average, ult = upstream ultimate).
\end{tablenotes}
\end{threeparttable}
\end{table*}


% Add after describing hydrological signatures:

A comprehensive set of hydrological signatures was calculated for all high-quality gauges 
(Table~\ref{tab:hydro_signatures}). These metrics characterize flow magnitude, timing, duration, 
baseflow dynamics, and water balance, providing process-based descriptors of catchment behavior.

\begin{table}[htbp]
\centering
\caption{Hydrological signatures computed for CAMELS-RU catchments. All signatures derived from daily discharge time series for gauges with $>$70\% data completeness per hydrological year and at least 5 qualifying years ($n = \nsigngauges$; half-flow date requires $>$82\% completeness, $n = \nhalfflowgauges$). FDC slope is positive for steeper curves (opposite sign convention to \citealt{Sawicz2011}).}
\label{tab:hydro_signatures}
\begin{tabular}{llp{5.5cm}}
\toprule
\textbf{Signature} & \textbf{Unit} & \textbf{Description} \\
\midrule
\multicolumn{3}{l}{\textit{Magnitude}} \\
\quad q\_mean & mm/day & Mean daily discharge \\
\quad runoff\_ratio & -- & Annual $Q/P$ ratio (ERA5-Land precipitation) \\
\midrule
\multicolumn{3}{l}{\textit{Variability}} \\
\quad q\_cv & -- & Coefficient of variation of daily discharge \\
\quad FDC\_slope & -- & FDC slope between 33rd and 66th exceedance percentiles (log-space; positive = steeper) \\
\quad flashiness\_index & -- & Richards--Baker index: $\sum|Q_{i}-Q_{i-1}|/\sum Q_i$ \\
\midrule
\multicolumn{3}{l}{\textit{Extremes}} \\
\quad Q5 & mm/day & 5th exceedance percentile (high flows) \\
\quad Q95 & mm/day & 95th exceedance percentile (low flows) \\
\quad high\_flow\_freq & \% of days & Days exceeding $2\times$ median discharge \\
\quad high\_flow\_dur & days & Mean duration of high-flow events ($>2\times$ median) \\
\midrule
\multicolumn{3}{l}{\textit{Baseflow}} \\
\quad baseflow\_index & -- & Baseflow fraction; ensemble of 1{,}000 Lyne--Hollick filter runs, $\alpha \sim U[0.9, 0.98]$, 3-pass \\
\midrule
\multicolumn{3}{l}{\textit{Low Flow}} \\
\quad low\_flow\_freq & \% of days & Days below $0.2\times$ mean discharge \\
\quad low\_flow\_dur & days & Mean duration of low-flow events ($<0.2\times$ mean) \\
\midrule
\multicolumn{3}{l}{\textit{Seasonality}} \\
\quad half\_flow\_date & day of hydro yr & Day of hydrological year (Oct~1 = day~1) when 50\% of annual streamflow volume has passed \\
\bottomrule
\end{tabular}
\end{table}


% In section 05_clustering_regional.tex:
% Add after describing clustering results:

Hierarchical clustering of HydroATLAS attributes identified \nclusters distinct catchment types 
(Table~\ref{tab:geo_clusters}). Clusters range from cropland-dominated lowlands (Cluster 1) to 
permafrost-affected highlands (Cluster 4), capturing the gradient of environmental conditions 
across Russia's vast territory.

\begin{table*}[htbp]
\centering
\caption{Geophysical clustering results: key characteristics of 15 catchment clusters based on HydroATLAS attributes.}
\label{tab:geo_clusters}
\small
\begin{tabular}{clrrrrrrr}
\toprule
\textbf{ID} & \textbf{Cluster Name} & \textbf{N} & \textbf{Area} & \textbf{Elev.} & \textbf{Precip.} & \textbf{Forest} & \textbf{Crop} & \textbf{Perm.} \\
 &  & \textbf{Gauges} & \textbf{(km$^2$)} & \textbf{(m)} & \textbf{(mm/yr)} & \textbf{(\%)} & \textbf{(\%)} & \textbf{(\%)} \\
\midrule
1 & Cropland / Clay-rich & 179 & 4659 & 202 & 532 & 3.0 & 64.0 & 0.0 \\
2 & Lowland / Moderate P & 238 & 4971 & 228 & 493 & 19.3 & 39.1 & 0.4 \\
3 & Silty / Snow-dominated & 59 & 3589 & 693 & 869 & 38.2 & 0.0 & 28.6 \\
4 & Permafrost / Snow & 110 & 11377 & 705 & 423 & 28.6 & 0.0 & 88.0 \\
5 & Forested / Sandy & 182 & 11104 & 1298 & 517 & 68.7 & 0.8 & 60.0 \\
6 & Forested / Permafrost & 105 & 9970 & 732 & 433 & 77.9 & 0.0 & 75.0 \\
7 & Snow-dominated & 46 & 10327 & 108 & 516 & 57.7 & 0.8 & 27.1 \\
8 & Sandy / Snow & 159 & 2701 & 159 & 561 & 59.6 & 2.7 & 2.0 \\
9 & Forested / Snow & 480 & 6570 & 321 & 605 & 88.8 & 2.1 & 5.6 \\
10 & Forested & 242 & 6336 & 382 & 539 & 57.4 & 17.7 & 6.4 \\
11 & Forested / High-PET & 67 & 3011 & 807 & 766 & 78.8 & 4.9 & 11.2 \\
12 & Forested / Karst & 253 & 5428 & 411 & 606 & 84.1 & 4.1 & 4.4 \\
13 & Karst / Clay-rich & 47 & 816 & 638 & 703 & 38.9 & 22.2 & 0.5 \\
14 & Protected / Karst & 21 & 498 & 1811 & 1023 & 34.6 & 3.8 & 11.0 \\
15 & Highland / Deep-GW & 63 & 7080 & 2215 & 813 & 16.5 & 10.8 & 19.7 \\
\midrule
\multicolumn{2}{l}{\textbf{Total / Mean}} & \textbf{2251} & \textbf{6080} & \textbf{513} & \textbf{574} & \textbf{52.1} & \textbf{14.5} & \textbf{21.3} \\
\bottomrule
\end{tabular}
\begin{tablenotes}
\item Clustering performed using Ward's method on 22 normalized HydroATLAS attributes (catchments $<$ 50\,000 km$^2$).
\item Area: median catchment area; Elev.: mean elevation; Precip.: mean annual precipitation; Forest/Crop/Perm.: percentage land cover.
\item Silhouette score: 0.194; Calinski--Harabasz index: 314.9. See Section~\ref{sec:physiographic_regionalization} for full analysis.
\end{tablenotes}
\end{table*}


% In section 06b_trend_analysis.tex:
% Add after discussing precipitation trends:

Long-term trend analysis reveals minimal temporal non-stationarity in precipitation across all 
three datasets (Table~\ref{tab:precipitation_trends}). Only 1.3--6.0\% of gauges show statistically 
significant trends ($p < 0.05$), with MSWEP exhibiting the highest percentage (6.0\%, all increasing). 
This temporal stability supports the use of the \years period for large-sample hydrology studies.

\begin{table*}[htbp]
\centering
\caption{Precipitation trend analysis for three datasets (2008--2023).}
\label{tab:precipitation_trends}
\begin{tabular}{lrrrrr}
\toprule
\textbf{Variable} & \textbf{N Gauges} & \textbf{Significant} & \textbf{Increasing} & \textbf{Decreasing} & \textbf{Mean Slope} \\
 & & \textbf{(\%)} & \textbf{(\%)} & \textbf{(\%)} & \textbf{(mm/yr$^2$)} \\
\midrule
Precipitation (ERA5-Land) & 1721 & 1.3 & 1.1 & 0.2 & 0.09 \\
Precipitation (MSWEP) & 1721 & 6.0 & 6.0 & 0.0 & 0.24 \\
Precipitation (GPCP) & 1721 & 1.6 & 1.0 & 0.6 & 0.25 \\
\bottomrule
\end{tabular}
\begin{tablenotes}
\small
\item Trends calculated using Mann-Kendall test with Sen's slope estimator; significance at $p < 0.05$.
\item Significant trends are minimal across all datasets ($<$6\%), indicating temporal stability during study period.
\item MSWEP shows higher percentage of significant trends but remains $<$10\% of gauges.
\end{tablenotes}
\end{table*}

